% Theseus — Investment memo (Round 19 prompt 11)
%
% Restrained typography: 11pt body, Palatino (with Charter as the
% fallback when the system has no Palatino), narrow margins, small
% cover page. The memo's PDF is the canonical artifact a portfolio
% agent or human reviewer reads.
%
% Token substitutions handled by noosphere/synthesizer/memo_pdf.py:
%   Bilateral defense spending grows >10\% over the next year under the cited…      — memo title (LaTeX-escaped)
%   Bilateral defense spending grows >10\% over the next year under the cited regime. Confidence 0.55–0.75. The governing principle is `principle\_security\_dilemma`.       — TL;DR paragraph (LaTeX-escaped, ≤ 80 words)
%   Portfolio Agent — investment  — who the memo is addressed to
%   2026-05-16       — YYYY-MM-DD
%   Theseus — synthesizer/v1     — "Theseus — synthesizer/vN"
%   \subsection*{Header}
\textbf{Title}: Bilateral defense spending grows >10\% over the next year under the cited…
\textbf{Author}: Theseus — synthesizer/v1
\textbf{Date}: 2026-05-16
\textbf{Question type}: INVESTMENT\_DECISION
\textbf{Confidence band}: 0.55–0.75
\textbf{Addressee}: Portfolio Agent — investment

\subsection*{TL;DR}
Bilateral defense spending grows >10\% over the next year under the cited regime. Confidence 0.55–0.75. The governing principle is \texttt{principle\_security\_dilemma}.

\subsection*{Question constituted}
should we long the polymarket arms-race escalation contract under the current regime?

\subsection*{Governing principles}
\begin{itemize}
\item \textbf{principle\_security\_dilemma} — States in mutual threat perception engage in security-dilemma escalation absent credible commitment devices. (\_Political Philosophy, Strategy\_; [details](/principles/principle\_security\_dilemma))
\item \textbf{principle\_domestic\_lockin} — Domestic political incentives reinforce external escalation once initiated; reversal requires elite cost. (\_Political Philosophy\_; [details](/principles/principle\_domestic\_lockin))
\item \textbf{principle\_second\_derivative} — Arms races are predicted by the second derivative of military spending, not the first. (\_Strategy, Economics\_; [details](/principles/principle\_second\_derivative))
\end{itemize}

\subsection*{Observed inputs}
\_No observed inputs recorded for this memo.\_

\subsection*{Reasoning chain}
\textbf{Step 1 — DETECT}: applied principle \texttt{principle\_security\_dilemma} to observation \texttt{—} → derived: Bilateral spending acceleration crosses the threshold.

\textbf{Step 2 — APPLY\_PRINCIPLE}: applied principle \texttt{principle\_domestic\_lockin} to observation \texttt{—} → derived: Domestic lock-in lowers the probability of reversal.

\textbf{Step 3 — SYNTHESIZE}: applied principle \texttt{principle\_security\_dilemma} to observation \texttt{—} → derived: Combine principles into a forward projection.

\subsection*{Implied bet}
\begin{itemize}
\item \textbf{Bet kind}: unspecified
\item \textbf{Side}: YES
\item \textbf{Stake}: 50.0
\item \textbf{Horizon}: —
\item \textbf{Ceilings}: —
\end{itemize}

Eight-gate readiness:
\begin{itemize}
\item ✅ \texttt{thesis\_articulated}
\item ✅ \texttt{principles\_govern}
\item ✅ \texttt{no\_standing\_contradiction}
\item ✅ \texttt{confidence\_band\_narrow}
\item ✅ \texttt{stake\_sized}
\item ⬜ \texttt{horizon\_set}
\item ⬜ \texttt{exit\_condition\_defined}
\item ✅ \texttt{addressee\_authorised}
\end{itemize}

\subsection*{What would update us}
We would weaken on any of: (1) principle \texttt{principle\_security\_dilemma} losing STANDING in the firm's contradiction lifecycle; (2) a fresh observation that flips the governing precondition; (3) a provenance downgrade on any cited principle. We would strengthen on a confirming algorithm invocation.

\subsection*{Abstentions and caveats}
Confidence band rationale: the synthesizer narrowed to a 0.20-wide band (low=0.55, high=0.75). No STANDING contradictions block the chain (else the synthesizer would have abstained).

\subsection*{Provenance audit}
Active provenance kinds:
\begin{itemize}
\item \textbf{PROPRIETARY} — weighting 2.00; sources: 0
\item \textbf{ENDORSED\_EXTERNAL} — weighting 1.00; sources: 0
\item \textbf{STUDIED\_EXTERNAL} — weighting 1.00; sources: 0
\end{itemize}       — rendered LaTeX body of the 10 sections
%
\documentclass[11pt]{article}

\usepackage[margin=1in]{geometry}
\usepackage{microtype}
\usepackage{titlesec}
\usepackage{enumitem}
\usepackage[hidelinks]{hyperref}
\usepackage{graphicx}
\usepackage[T1]{fontenc}
\usepackage{textcomp}

% Charter is the founder preference; fall back gracefully when Charter
% is not installed (CI containers, minimal LaTeX distros).
\IfFileExists{charter.sty}{\usepackage{charter}}{%
  \IfFileExists{mathpazo.sty}{\usepackage{mathpazo}}{}%
}

\setlist[itemize]{leftmargin=*, topsep=2pt, itemsep=2pt}
\setlist[enumerate]{leftmargin=*, topsep=2pt, itemsep=2pt}
\titleformat{\section}{\large\bfseries}{}{0pt}{}
\titlespacing{\section}{0pt}{14pt}{4pt}
\titleformat{\subsection}{\normalsize\bfseries}{}{0pt}{}
\titlespacing{\subsection}{0pt}{10pt}{2pt}

\pagestyle{plain}
\setlength{\parskip}{0.6em}
\setlength{\parindent}{0pt}

\title{Bilateral defense spending grows >10\% over the next year under the cited…}
\author{Theseus — synthesizer/v1}
\date{2026-05-16}

\begin{document}

% ── Cover page ─────────────────────────────────────────────────────
\thispagestyle{empty}
\begin{flushleft}
  {\footnotesize\textsc{Theseus} \quad \textsc{Investment memo}}\\[4pt]
  \rule{\linewidth}{0.4pt}\\[18pt]
  {\LARGE\bfseries Bilateral defense spending grows >10\% over the next year under the cited…}\\[12pt]
  {\small\itshape Addressed to: Portfolio Agent — investment}\\[6pt]
  {\small Date: 2026-05-16 \quad Author: Theseus — synthesizer/v1}\\[24pt]
  \rule{\linewidth}{0.4pt}\\[14pt]
  {\small\textbf{TL;DR.} Bilateral defense spending grows >10\% over the next year under the cited regime. Confidence 0.55–0.75. The governing principle is `principle\_security\_dilemma`.}
\end{flushleft}
\vfill
\begin{center}
  \footnotesize\textsc{Confidential — for the portfolio agent named above. Do not redistribute without operator approval.}
\end{center}
\newpage

% ── Body ────────────────────────────────────────────────────────────
\subsection*{Header}
\textbf{Title}: Bilateral defense spending grows >10\% over the next year under the cited…
\textbf{Author}: Theseus — synthesizer/v1
\textbf{Date}: 2026-05-16
\textbf{Question type}: INVESTMENT\_DECISION
\textbf{Confidence band}: 0.55–0.75
\textbf{Addressee}: Portfolio Agent — investment

\subsection*{TL;DR}
Bilateral defense spending grows >10\% over the next year under the cited regime. Confidence 0.55–0.75. The governing principle is \texttt{principle\_security\_dilemma}.

\subsection*{Question constituted}
should we long the polymarket arms-race escalation contract under the current regime?

\subsection*{Governing principles}
\begin{itemize}
\item \textbf{principle\_security\_dilemma} — States in mutual threat perception engage in security-dilemma escalation absent credible commitment devices. (\_Political Philosophy, Strategy\_; [details](/principles/principle\_security\_dilemma))
\item \textbf{principle\_domestic\_lockin} — Domestic political incentives reinforce external escalation once initiated; reversal requires elite cost. (\_Political Philosophy\_; [details](/principles/principle\_domestic\_lockin))
\item \textbf{principle\_second\_derivative} — Arms races are predicted by the second derivative of military spending, not the first. (\_Strategy, Economics\_; [details](/principles/principle\_second\_derivative))
\end{itemize}

\subsection*{Observed inputs}
\_No observed inputs recorded for this memo.\_

\subsection*{Reasoning chain}
\textbf{Step 1 — DETECT}: applied principle \texttt{principle\_security\_dilemma} to observation \texttt{—} → derived: Bilateral spending acceleration crosses the threshold.

\textbf{Step 2 — APPLY\_PRINCIPLE}: applied principle \texttt{principle\_domestic\_lockin} to observation \texttt{—} → derived: Domestic lock-in lowers the probability of reversal.

\textbf{Step 3 — SYNTHESIZE}: applied principle \texttt{principle\_security\_dilemma} to observation \texttt{—} → derived: Combine principles into a forward projection.

\subsection*{Implied bet}
\begin{itemize}
\item \textbf{Bet kind}: unspecified
\item \textbf{Side}: YES
\item \textbf{Stake}: 50.0
\item \textbf{Horizon}: —
\item \textbf{Ceilings}: —
\end{itemize}

Eight-gate readiness:
\begin{itemize}
\item ✅ \texttt{thesis\_articulated}
\item ✅ \texttt{principles\_govern}
\item ✅ \texttt{no\_standing\_contradiction}
\item ✅ \texttt{confidence\_band\_narrow}
\item ✅ \texttt{stake\_sized}
\item ⬜ \texttt{horizon\_set}
\item ⬜ \texttt{exit\_condition\_defined}
\item ✅ \texttt{addressee\_authorised}
\end{itemize}

\subsection*{What would update us}
We would weaken on any of: (1) principle \texttt{principle\_security\_dilemma} losing STANDING in the firm's contradiction lifecycle; (2) a fresh observation that flips the governing precondition; (3) a provenance downgrade on any cited principle. We would strengthen on a confirming algorithm invocation.

\subsection*{Abstentions and caveats}
Confidence band rationale: the synthesizer narrowed to a 0.20-wide band (low=0.55, high=0.75). No STANDING contradictions block the chain (else the synthesizer would have abstained).

\subsection*{Provenance audit}
Active provenance kinds:
\begin{itemize}
\item \textbf{PROPRIETARY} — weighting 2.00; sources: 0
\item \textbf{ENDORSED\_EXTERNAL} — weighting 1.00; sources: 0
\item \textbf{STUDIED\_EXTERNAL} — weighting 1.00; sources: 0
\end{itemize}

\end{document}
